\section*{Nudelauflauf mit Rahmspinat und Meatballs}
\ihead{}\chead{Personen:4}\ohead{}
\ifoot{Vorbereitungszeit:30 min}\cfoot{}\ofoot{Kochzeit:40 min}
{\Large Zutaten}
\begin{multicols}{2}
\begin{itemize}
\item \SI{500}{g} grüne, frische Tagliatelle
\item \num{1} Zwiebel
\item \num{1} Knoblauchzehe
\item \SI{500}{g} Rahmspinat, TK
\item \SI{100}{g} Frischkäse
\item \SI{100}{g} Sahne
\item \SI{200}{g} italienische Bratwurst, roh
\item \SI{100}{g} junger Pecorino, frisch gerieben
\item Olivenöl
\item Salz und Pfeffer aus der Mühle
\item frisch geriebene Muskatnuss
\end{itemize}
% \columnbreak
\end{multicols}
\noindent {\Large Zubereitung}\\
\\
Die Tagliatelle entsprechend der Anleitung auf der Packung zubereiten und zur Seite stellen.
Die Zwiebel in einem Topf mit Olivenöl anschwitzen, bis sie glasig ist.
Den Knoblauch dazupressen und sofort den TK-Spinat dazugeben.
Warten, bis der Spinat aufgetaut ist, und dabei immer wieder umrühren, damit nichts anbrennt.
Frischkäse und Sahne dazugeben und weiterrühren, bis der Frischkäse geschmolzen ist.
Mit Salz, Pfeffer und Muskat abschmecken.
Nudeln und Spinatsoße in einer Auflaufform miteinander mischen.
\textit{Vorbereitung:} Den Ofen auf \SI{200}{\celsius} Ober-/Unterhitze vorheizen.
Die italienische Bratwurst in Stücke schneiden und mit den Händen zu kleinen Bällchen formen.
Die Fleischbällchen in einer Pfanne mit Olivenöl schön knusprig braten und anschließend auf der Nudel-Spinat-Mischung verteilen.
Den Pecorino darüberstreuen und nach Belieben noch etwas Salz, Pfeffer und Muskat darübergeben.
Den Auflauf für \SI{20-30}{min} in den heißen Backofen stellen, bis der Käse geschmolzen und goldgelb ist.
